\documentclass{article}
\usepackage[utf8]{inputenc}
\usepackage[english]{babel}
\usepackage{amsmath, amssymb}
\usepackage{amsthm}
\usepackage{mathrsfs}
\usepackage{mathtools}
\usepackage{subcaption}
\usepackage{graphicx}
\usepackage[colorlinks=true,      % 启用彩色链接而非带框链接
            linkcolor=blue,       % 内部链接（如目录、交叉引用）颜色为蓝色
            urlcolor=blue,        % 外部链接颜色为蓝色
            citecolor=blue]{hyperref} % 参考文献链接颜色为蓝色
\usepackage{indentfirst}
\usepackage{geometry}
\geometry{a4paper, margin=1.5cm}

\newtheorem{theorem}{Theorem}[section]
\newtheorem{lemma}[theorem]{Lemma}
\newtheorem{corollary}[theorem]{Corollary}
\newtheorem*{remark}{Remark}
\newtheorem{proposition}[theorem]{Proposition}

\theoremstyle{definition}
\newtheorem{definition}[theorem]{Definition}
\newtheorem{example}[theorem]{Example}
\newtheorem{problem}{Problem}
\newtheorem{solution}{Solution}

\newcommand{\R}{\mathbb{R}}
\newcommand{\N}{\mathbb{N}}
\newcommand{\Q}{\mathbb{Q}}
\newcommand{\C}{\mathbb{C}}
\newcommand{\Z}{\mathbb{Z}}
\newcommand{\dd}{\mathrm{d}}
\newcommand{\re}{\mathrm{Re}}
\newcommand{\im}{\mathrm{Im}}
\newcommand{\res}{\mathrm{Res}}
\newcommand{\Arg}{\mathrm{Arg}}
\newcommand{\hol}{\mathrm{Hol}}
\newcommand{\auth}{\mathrm{Aut}_{\mathrm{hol}}}
\renewcommand{\leq}{\leqslant}
\renewcommand{\geq}{\geqslant}
\renewcommand{\emptyset}{\varnothing}
\renewcommand{\le}{\leqslant}
\renewcommand{\ge}{\geqslant}

\title{\textbf{A Summary of Complex Analysis}}
\author{LUO Yunjie}
\date{\today}

\begin{document}
\maketitle
\begin{abstract}
    This document provides a summary of complex analysis. The main references are \emph{Complex Analysis} by Elias M. Stein and Rami Shakarchi and the lecture notes of \'Ecole Polytechnique.
\end{abstract}
\tableofcontents
\newpage
\section{Holomorphic Functions}
\subsection{Basic Definitions}
What is a holomorphic function? It is a complex function that is complex differentiable in a neighborhood of every point. Moreover, complex differentiability is equivalent to real differentiability plus the Cauchy-Riemann equations.
\begin{theorem}[Cauchy-Riemann Equations]
Let \(f\colon\Omega\to\C\) be a complex function defined on an open set \(\Omega\subset\C\). Write \(f(z)=u(x,y)+i v(x,y)\) where \(z=x+i y\), and \(u,v\colon\Omega^{\prime}\subset\R^2\to\R\) are real-valued functions. 
Then \(f\) is complex differentiable at \(z_0=x_0+iy_0\) if and only if \(u\) and \(v\) are real differentiable at \((x_0,y_0)\) and satisfy the Cauchy-Riemann equations at \((x_0,y_0)\):
\[\frac{\partial u}{\partial x} = \frac{\partial v}{\partial y}, \quad \frac{\partial u}{\partial y} = -\frac{\partial v}{\partial x}.\]
\end{theorem}
By the Cauchy-Riemann equations, if \( f \) is holomorphic on a connected open set, then \( \re(f) \) determines \( \im(f) \) up to an additive constant, and conversely.

In addition, we can also consider 
\begin{equation*}
    \frac{\partial }{\partial z}\coloneq \frac{1}{2}\left(\frac{\partial }{\partial x} - i \frac{\partial }{\partial y}\right), \quad \frac{\partial }{\partial \bar{z}} \coloneq \frac{1}{2}\left(\frac{\partial }{\partial x} + i \frac{\partial }{\partial y}\right).
\end{equation*}
Then the Cauchy-Riemann equations are equivalent to \(\frac{\partial f}{\partial \bar{z}}=0\). Here are some properties about $\frac{\partial }{\partial z}$ and $\frac{\partial }{\partial \bar{z}}$.
\begin{proposition}
If \(f\) and \(g\) are complex-valued functions, then
\begin{itemize}
    \item $\overline{\frac{\partial f}{\partial z}}=\frac{\partial \bar{f}}{\partial \bar{z}}$.
    \item $\frac{\partial (f\circ g)}{\partial z}(z)
    =\left(\frac{\partial f}{\partial z}\circ g\right)(z)\cdot \frac{\partial g}{\partial z}(z)+\left(\frac{\partial f}{\partial \bar{z}}\circ g\right)(z)\cdot \frac{\partial \bar{g}}{\partial z}(z).$
\end{itemize}
\end{proposition}
Using these operators, we can easily verify that $\overline{f(\bar{z})}$ is holomorphic if \(f\) is holomorphic. 

\subsection{Integration along Curves}\label{integration along curves}
We now turn to an important property of holomorphic functions. A holomorphic function is infinitely complex differentiable. To prove this, we first introduce integration along curves.

Given a smooth curve \(\gamma\colon[a,b]\to\C\) and a complex function \(f\colon\Omega\to\C\) defined on an open set \(\Omega\subset\C\) containing the image of \(\gamma\), we define the complex integral of \(f\) along \(\gamma\) by
\begin{equation*}
    \int_{\gamma} f(z)\, \dd z = \int_a^b f(\gamma(t)) \gamma'(t)\, \dd t.
\end{equation*}
\begin{proposition}\label{homotopic prop}
Suppose $\Omega$ is a connected open set. If a continuous function \(f\colon\Omega\to\C\) has a primitive $F$ in \(\Omega\), and \(\gamma\colon[a,b]\to\Omega\) is a smooth curve in \(\Omega\), then
\begin{equation*}
    \int_{\gamma} f(z)\, \dd z = F(\gamma(b)) - F(\gamma(a)).
\end{equation*}
\end{proposition}
With the definition above, we can state Cauchy's integral formulas.
\begin{theorem}[Cauchy's Integral Formulas]\label{th:cauchy_integral_formula}
Let \(\Omega\subset\C\) be an open set containing $\overline{D(z_0,r)}$, and let \(f\colon\Omega\to\C\) be a holomorphic function. Then
\begin{equation*}
    f(z) = \frac{1}{2\pi i} \int_{C(z_0,r)} \frac{f(\omega)}{\omega-z}\, \dd \omega,
\end{equation*}
for all $z\in D(z_0,r)$.
\end{theorem}
The key idea is to consider \[\frac{f(\omega)}{\omega - z}=\frac{f(\omega)-f(z)}{\omega - z}+\frac{f(z)}{\omega - z} \]
and the keyhole contour.

Using Cauchy's integral formulas, we can easily prove that holomorphic functions are infinitely complex differentiable. 
\begin{corollary}
    Let \(\Omega\subset\C\) be an open set containing $\overline{D(z_0,r)}$, and let \(f\colon\Omega\to\C\) be a holomorphic function. Then
\begin{equation*}
    f^{(n)}(z) = \frac{n!}{2\pi i} \int_{C(z_0,r)} \frac{f(\omega)}{(\omega-z)^{n+1}}\, \dd \omega,
\end{equation*}
for all $z\in D(z_0,r)$.
\end{corollary}
Thus holomorphic functions are analytic, meaning they have power series expansions. Conversely, any convergent power series defines a holomorphic function, since term-by-term differentiation is valid.
We can write $f(z)=\sum_{n\geq 0} a_n (z-z_0)^n$ in $D(z_0,r)\subseteq \Omega$ with 
\begin{align*}
a_n&=\frac{f^{(n)}(z_0)}{n!}=\frac{1}{2\pi i} \int_{C(z_0,r)} \frac{f(\omega)}{(\omega-z_0)^{n+1}}\, 
\dd \omega\\
&=\frac{1}{2\pi r^n}\int_{0}^{2\pi}f\left(z_0+re^{i\theta}\right)e^{-in\theta}\, \dd\theta.
\end{align*}
Denote by $\hol(\Omega)$ the set of holomorphic functions on $\Omega$ and $\mathcal{A}(\Omega)$ the set of analytic functions on $\Omega$. Since $\hol(\Omega)=\mathcal{A}(\Omega)$, we can conclude the existence of a primitive of holomorphic functions on a disc.
\begin{proposition}
Let $\Omega\subset\C$ be an open set, and let $\gamma_1,\gamma_2$ be two piecewise smooth curves in $\Omega$ that are homotopic in $\Omega$. Let \(f\colon\Omega\to\C\) be a holomorphic function, then
\begin{equation*}
    \int_{\gamma_1} f(z)\, \dd z = \int_{\gamma_2} f(z)\, \dd z.
\end{equation*}
\end{proposition}
The proof is nontrivial. There are many details. The key idea is to cover two nearby curves with disks since holomorphic functions have primitives on disks.

As a consequence of Cauchy's estimates, we obtain Liouville's theorem.
\begin{theorem}[Liouville's Theorem]
Let \(f\colon\C\to\C\) be a bounded entire function. Then \(f\) is constant.
\end{theorem}
\begin{remark}
    Liouville's theorem yields a short proof of the fundamental theorem of algebra.
\end{remark}
\subsection{Analytic Continuation and Uniqueness Theorem}
We now turn to analytic continuation. The uniqueness of analytic continuation follows from the fact that the zeros of a non-zero holomorphic function are isolated.
\begin{theorem}[Isolation of Zeros]
Let \(\Omega\subset\C\) be a connected open set, and let \(f\colon\Omega\to\C\) be a holomorphic function. If there exists a point \(z_0\in\Omega\) such that \(f(z_0)=0\) and \(f\) is not identically zero in any neighborhood of \(z_0\), then there exists a neighborhood \(D(z_0,r)\subseteq \Omega\) such that \(f(z)\neq 0\) for all \(z\in D(z_0,r)\setminus\{z_0\}\).
\end{theorem}
The key idea of the proof is to consider the set $U\coloneq \{ z\in \Omega\colon f^{(n)}(z)=0,\,\forall n\geq 0\}$, which is closed (since it is an intersection of closed sets) and open (since the Taylor series shows that it is identically zero in a neighborhood). 

Then we obtain the identity theorem which is proved by the isolation of zeros of holomorphic functions.
\begin{theorem}[Identity Theorem]
Let \(\Omega\subset\C\) be a connected open set, and let \(f,g\colon\Omega\to\C\) be holomorphic functions. If there exists a subset \(S\subset\Omega\) with an accumulation point in \(\Omega\) such that \(f(z)=g(z)\) for all \(z\in S\), then \(f(z)=g(z)\) for all \(z\in\Omega\).
\end{theorem}
    This theorem implies the uniqueness of analytic continuation.

\subsection{Sequences of Holomorphic Functions}
Denote $\|f\|_{\infty,U}=\sup\limits_{z\in U}|f(z)|$ for a function $f\colon U\to \C$ and a set $U\subseteq \C$.
\begin{definition}
    Let \(\Omega\subset\C\) be an open set, and let \((f_n)\) be a sequence of complex-valued functions defined on \(\Omega\).
    We say that \((f_n)\) converges locally uniformly if there exists a function $f\colon \Omega \to \C$
    such that for all $z\in\Omega$, there exists a neighborhood $U_z\subseteq \Omega$ such that \[\lim_{n\to \infty} \|f_n - f\|_{\infty,U_z} = 0.\]
    We say that \((f_n)\) is a local Cauchy sequence if for all $z\in\Omega$, there exists a neighborhood $U_z\subseteq \Omega$ such that 
    \((f_n)\) is a Cauchy sequence with respect to the norm \(\|\cdot\|_{\infty,U_z}\), i.e., for all \(\varepsilon>0\), there exists \(N>0\) such that for all \(m,n>N\), 
    \[\|f_n - f_m\|_{\infty,U_z} < \varepsilon.\]
\end{definition}
\begin{proposition}\label{prop:seq equiv}
    Let \(\Omega\subset\C\) be an open set, and let \((f_n)\) be a sequence of complex-valued functions defined on \(\Omega\). 
    Then \((f_n)\) converges locally uniformly to a function \(f\colon\Omega\to\C\) if and only if for every compact subset \(K\subseteq \Omega\), \((f_n)\) converges uniformly to \(f\) on \(K\).
    Moreover, \((f_n)\) is a local Cauchy sequence if and only if for every compact subset \(K\subseteq \Omega\), \((f_n)\) is a Cauchy sequence with respect to the norm \(\|\cdot\|_{\infty,K}\).
\end{proposition}
The proof is simple by using the compactness of $K$ and the fact that $\C$ is locally compact.
\begin{proposition}
    Let \(\Omega\subset\C\) be an open set, and let \((f_n)\) be a sequence of complex-valued functions defined on \(\Omega\). Then TFAE:
    \begin{itemize}
        \item \((f_n)\) converges locally uniformly.
        \item \((f_n)\) is a local Cauchy sequence.
    \end{itemize}
\end{proposition}
This follows immediately from Proposition \ref{prop:seq equiv}.
\begin{theorem}[Weierstrass Convergence Theorem]
    Let \(\Omega\subset\C\) be an open set, and let \((f_n)\) be a sequence of holomorphic functions defined on \(\Omega\). If \((f_n)\) converges locally uniformly to a function \(f\colon\Omega\to\C\), then 
    \begin{itemize}
        \item \(f\) is holomorphic on \(\Omega\).
        \item For all \(k\geq 0\), the sequence \(\left(f_n^{(k)}\right)\) converges locally uniformly to \(f^{(k)}\) on \(\Omega\).
    \end{itemize}
\end{theorem}
\begin{proof}
    Given a point \(z_0\in\Omega\), choose \(r>0\) such that \(\overline{D(z_0,r)}\subseteq \Omega\). Since the $f_n$ are continuous and converge uniformly to $f$, $f$ is continuous. For all \(z\in D(z_0,r)\), we can define
    \begin{equation*}
        F(z) = \frac{1}{2\pi i} \int_{C(z_0,r)} \frac{f(\omega)}{\omega-z}\, \dd \omega.
    \end{equation*}
    $F$ is holomorphic on $D(z_0,r)$ since $f$ is continuous. Since \((f_n)\) converges locally uniformly to \(f\) on \(\Omega\), by Proposition \ref{prop:seq equiv} it converges uniformly to \(f\) on \(\overline{D(z_0,r)}\). 
    Hence for all \(\varepsilon>0\), there exists \(N>0\) such that for all \(n>N\) and all \(\omega\in C(z_0,r)\), \(|f_n(\omega)-f(\omega)|<\varepsilon\). Therefore, for all \(z\in D(z_0,r/2)\) and all \(n>N\),
    \begin{align*}
        \left|f_n(z)-\frac{1}{2\pi i} \int_{C(z_0,r)} \frac{f(\omega)}{\omega-z}\, \dd \omega\right| & = \left|\frac{1}{2\pi i} \int_{C(z_0,r)} \frac{f_n(\omega)-f(\omega)}{\omega-z}\, \dd \omega\right|\\
        & \leq \frac{1}{2\pi} \int_{C(z_0,r)} \frac{|f_n(\omega)-f(\omega)|}{|\omega-z|}\, \dd \omega\\
        & < \frac{\varepsilon}{2\pi} \int_{C(z_0,r)} \frac{1}{|\omega-z|}\, \dd \omega\\
        & \leq  \frac{\varepsilon\cdot r}{r-|z-z_0|}\\
        &\leq 2\varepsilon.
    \end{align*}
    Hence, \(f(z)=F(z)\) for all \(z\in D(z_0,r/2)\). We conclude that \(f\) is holomorphic on \(\Omega\). The second part can be proved similarly by using the Cauchy Integral Formulas for derivatives.
\end{proof}
\begin{corollary}
    If $f=\sum f_n$ is a series of holomorphic functions converging locally uniformly on an open set \(\Omega\subset\C\), then \(f\) is holomorphic on \(\Omega\) and for all \(k\geq 0\), 
    \begin{equation*}
        f^{(k)} = \sum f_n^{(k)},
    \end{equation*}
    where the series on the right converges locally uniformly on \(\Omega\).
\end{corollary}


\subsection{Holomorphic Functions Defined in Terms of Integrals}
Finally, we consider holomorphic functions defined by integrals. Suppose $\Omega$ is an open set and $X$ is a measurable subset of $\R^N$. 
If $F\colon\Omega\times X\to\C$ is a continuous function, then we consider the function defined by
\begin{equation*}
    f(z) = \int_X F(z,x)\, \dd x.
\end{equation*}
 
\begin{theorem}\label{thm:holomorphic-integral}
    Suppose that \(f\) is well-defined for all \(z\in \Omega\), and that for each \(x\in X\), the function \(F(\cdot,x)\colon\Omega\to\mathbb C\) is holomorphic. Then \(f\) is holomorphic on \(\Omega\) if either of the following conditions holds:
    \begin{enumerate}
        \item \(X\) is compact (and \(F\) is continuous on \(\Omega\times X\));
        \item for each \(z_0\in \Omega\), there exist a neighborhood \(D(z_0,r_0)\subseteq \Omega\) and an integrable function \(g_{z_0}\colon X\to[0,\infty)\) such that
        \[
            |F(z,x)| \leq g_{z_0}(x)
        \]
        for all \(z\in D(z_0,r_0)\) and almost all \(x\in X\).
    \end{enumerate}
    Moreover, under either condition, for every \(k\ge 1\) and every \(z\in \Omega\),
    \begin{equation*}
        f^{(k)}(z) = \int_X \frac{\partial^k F}{\partial z^k}(z,x)\, \dd x.
    \end{equation*}
\end{theorem}
\begin{remark}
    The local domination condition on $F$ ensures that we may interchange integration with differentiation.
\end{remark}
\section{Meromorphic Functions and Residue Theorem}
\subsection{Basic Definitions}
\begin{definition}
    $P$ is a subset of $\C$.
    \begin{itemize}
        \item $P$ is called discrete if all points in $P$ are isolated, i.e., for all $z\in P$, there exists $r>0$ such that $D(z,r)\cap P=\{z\}$.
        \item Suppose $\Omega$ is an open set in $\C$ containing $P$. $P$ is called locally finite in $\Omega$ if for all $a\in \Omega$, there exists a neighborhood $U_a$ such that $P\cap U_a$ is finite, or equivalently, for every compact subset $K\subseteq \Omega$, $P\cap K$ is finite.
    \end{itemize}
\end{definition} 
\begin{proposition}
    Suppose $\Omega$ is an open set in $\C$ and $P$ is a closed subset of $\Omega$. TFAE: 
    \begin{itemize}
        \item $P$ is discrete.
        \item $P$ is locally finite in $\Omega$.
        \item $P$ doesn't have accumulation points in $\Omega$. ($a$ is an accumulation point of $P$ if for all $r>0$, $(D(a,r)\cap P)\setminus\{a\}\neq \varnothing$.)
    \end{itemize}
\end{proposition}
\begin{definition}
    Let \(\Omega\subset\C\) be an open set, and let $z_0$ be a point in $\Omega$, and let \(f\colon\Omega\setminus\{z_0\}\to\C\) be holomorphic. 
    \begin{itemize}
        \item If there exists a holomorphic function \(g\colon\Omega\to\C\) such that \(g(z)=f(z)\) for all \(z\in \Omega\setminus\{z_0\}\), then we say that \(f\) has a removable singularity at \(z_0\).
        \item If there exists an integer \(m\geq 1\) such that $(z-z_0)^m f(z)$ is holomorphic at $z_0$ and nonzero at $z_0$, then we say that \(f\) has a pole of order \(m\) at \(z_0\).
        \item If \(f\) has neither a removable singularity nor a pole at \(z_0\), then we say that \(f\) has an essential singularity at \(z_0\).
    \end{itemize}
\end{definition}
\begin{remark}
    \begin{enumerate}
        \item One can distinguish them by examining the behavior of $f$ near $z_0$. If $\lim_{z\to z_0} f(z)$ exists, then it is a removable singularity. If $\lim_{z\to z_0} |f(z)| = +\infty$, then it is a pole. Otherwise, it is an essential singularity.   
        \item An equivalent definition of a pole is as follows: \(f\) has a pole of order \(m\) at $z_0$ if and only if $1/f$ is a holomorphic function in a neighborhood of $z_0$ and has a zero of order \(m\) at $z_0$.
    \end{enumerate}
\end{remark}

\begin{definition}\label{meromorphic}
    Let \(\Omega\subset\C\) be an open set, and let \(P\) be a locally finite subset of \(\Omega\). A function \(f\colon\Omega\setminus P\to\C\) is called meromorphic on \(\Omega\) if \(f\) is holomorphic on \(\Omega\setminus P\) and each point of \(P\) is either a pole or a removable singularity.
\end{definition}
\begin{remark}
    Here is another definition of meromorphic functions: A function \(f\) is called meromorphic on \(\Omega\) if for every point $z$ of $\Omega$ there is a neighborhood of $z$ in which at least one of $f$ and $1/f$ is holomorphic. 
\end{remark}
Denote by $Z(f)$ the set of zeros of $f$ and $P(f)$ the set of poles of $f$. By Definition \ref{meromorphic}, we can easily verify that if $f$ is a nonzero meromorphic function on a connected open set $\Omega$, then $Z(f)$ and $P(f)$ are both locally finite in $\Omega$. 

\subsection{Residue Theorem}
\begin{definition}[Laurent Series and Residue]
    Let \(f\colon\Omega\to\C\) be a meromorphic function defined on an open set \(\Omega\subset\C\) and let $z_0$ be a point in \(\Omega\). Then there exists $m\in \N$ and $r>0$ such that 
    \begin{equation*}
        f(z) = \sum_{n=-m}^{+\infty} a_n (z-z_0)^n,
    \end{equation*}
    for all \(z\in D(z_0,r)\setminus\{z_0\}\). This series is called the Laurent series of \(f\) centered at \(z_0\). This series is unique and absolutely convergent on an annulus \(\{r_1<|z-z_0|<r_2\}\) with \(0\leq r_1<r_2<r\). 
    And $f$ has a pole of order $m$ at $z_0$ if and only if $m\geq 1$ and $a_{-m}\neq 0$. In particular,
    the coefficient \(a_{-1}\) is called the residue of \(f\) at \(z_0\), denoted by \(\res(f,z_0)\). Moreover, the coefficients \(a_n\) can be computed by
    \begin{equation*}
        a_n = \frac{1}{2\pi i} \int_{C(z_0,\rho)} \frac{f(\omega)}{(\omega-z_0)^{n+1}}\, \dd \omega,
    \end{equation*}
    for all \(n\geq -m\) and all \(\rho\in(r_1,r_2)\).
\end{definition}

\begin{example}
    We give some examples of computing residues.
    \begin{enumerate}
        \item If \(f\) has a removable singularity at \(z_0\), then \(\res(f,z_0)=0\).
        \item If \(f\) has a pole of order \(n\) at \(z_0\), then \[\res(f,z_0)=\frac{1}{(n-1)!}\lim_{z\to z_0}\frac{d^{n-1}}{dz^{n-1}}((z-z_0)^n f(z)).\]
        \item If we can write \(f(z)=\sum_{n=-m}^{+\infty} a_n (z-z_0)^n\) in a neighborhood of \(z_0\), then \(\res(f,z_0)=a_{-1}\), for example, for $f(z)=1/(z^2+1)$ we can write
        \begin{equation*}
            f(z)=\frac{1}{(z-i)(z+i)}=\frac{1}{2i}\cdot\frac{1}{z-i}-\frac{1}{2i}\cdot\frac{1}{z+i},
        \end{equation*}
        so we have \(\res(f,i)=\frac{1}{2i}\) and \(\res(f,-i)=-\frac{1}{2i}\).
        \item If \(f(z)=g(z)/{h(z)}\) where \(g\) and \(h\) are holomorphic functions in a neighborhood of \(z_0\), \(g(z_0)\neq 0\) and \(h\) has a simple zero at \(z_0\), then \(f\) has a pole of order \(1\) at \(z_0\) and 
        \begin{equation*}
            \res(f,z_0)=\frac{g(z_0)}{h^{\prime}(z_0)}.
        \end{equation*}
    \end{enumerate}
\end{example}
What about poles and residues at infinity? We can consider $g(w)=f(1/w)$. If $g$ has a pole of order $m$ at $0$, then we say that $f$ has a pole of order $m$ at infinity. And we define the residue of $f$ at infinity by 
\[\res(f,\infty) = -\res\left(\frac{1}{w^2}\cdot f\left(\frac{1}{w}\right),0\right).\]
If $g$ has an expansion in Laurent series at $0$ like $g(w)=\sum_{n=-m}^{+\infty} a_n w^n$, then we have 
\[\res(f,\infty) = -a_{1}.\]

Now we state the residue theorem.
\begin{theorem}[Residue Theorem]
Let \(\Omega\subset\C\) be a simply connected open set. Suppose \(f\) is a meromorphic function on $\Omega$. If \(\gamma\) is a closed smooth curve in \(\Omega\setminus P(f)\), then
\begin{equation*}
    \int_{\gamma} f(z)\, \dd z = 2\pi i \sum_{z_0\in P(f)} \res(f,z_0) \cdot \mathrm{Ind}(\gamma,z_0),
\end{equation*}
where \(\mathrm{Ind}(\gamma,z_0)\) is the winding number of \(\gamma\) around \(z_0\), i.e.,
\begin{equation*}
    \mathrm{Ind}(\gamma,z_0) = \frac{1}{2\pi i} \int_{\gamma} \frac{1}{z-z_0}\, \dd z=\frac{1}{2\pi i}\int_{a}^{b}\frac{\gamma^{\prime}(t)}{\gamma(t)-z_0}\, \dd t.
\end{equation*}
\end{theorem}
\begin{remark}
    The definition of winding number shows that it is always an integer. In particular, if \(\gamma\) is a positively oriented simple closed curve whose interior is contained in \(\Omega\), then \(\mathrm{Ind}(\gamma,z_0)=1\) for \(z_0\) in the interior of \(\gamma\), and \(\mathrm{Ind}(\gamma,z_0)=0\) otherwise.
\end{remark}
\begin{example}
    We give some examples of applying the residue theorem to compute real integrals. 
    \begin{enumerate}
        \item Compute \(\int_{C(0,1)} \frac{1}{z}\left(z+\frac{1}{z}\right)^n\, \dd z\) for \(n\in\N\). And deduce the value of \(\int_0^{2\pi} \cos^n \theta\, \dd \theta\).
            \begin{equation*}
                \int_{C(0,1)} \frac{1}{z}\left(z+\frac{1}{z}\right)^n\, \dd z = \begin{cases}
                    2\pi i \cdot \binom{n}{n/2}, & n\text{ is even},\\
                    0, & n\text{ is odd}.
                \end{cases}
            \end{equation*}
            \begin{equation*}
                \int_0^{2\pi} \cos^n \theta\, \dd \theta = \begin{cases}
                    \frac{\pi}{2^{n-1}}\cdot \binom{n}{n/2}, & n\text{ is even},\\
                    0, & n\text{ is odd}.
                \end{cases}
            \end{equation*}
        \item Compute $I_{m,n}=\int_{0}^{+\infty}\frac{x^n}{1+x^m}\, \dd x$ for $m,n\in \N$ and $m\geq n+2$.\\
        We consider the contour integral of $f(z)=\frac{z^n}{1+z^m}$ on a sector with angle $\frac{2\pi}{m}$.
            \begin{equation*}
                I_{m,n} = \frac{\pi}{m\sin\left(\frac{(n+1)\pi}{m}\right)}.
            \end{equation*}
        \item Compute $\int_{0}^{2\pi} \frac{\dd t}{a^2\cos^2 t+b^2\sin^2 t}$ for $a>b>0$.\\
        We consider the change of variables $z=e^{it}$.
            \begin{equation*}
                    \int_{0}^{2\pi} \frac{\dd t}{a^2\cos^2 t+b^2\sin^2 t} = \frac{2\pi}{ab}.
            \end{equation*}
        \item Compute $I(s)=\int_{0}^{+\infty}\frac{x^{s-1}}{1+x}\,\dd x$ for $0<\re(s)<1$.\\
        We consider the change of variables $x=e^{t}$ and the contour integral of $f(z)=\frac{e^{sz}}{1+e^z}$ on the rectangle $(-R,R,R+2\pi i,-R+2\pi i)$.
            \begin{equation*}
                        I(s) = \frac{\pi}{\sin(\pi s)}.
            \end{equation*}
    \end{enumerate}
\end{example}

\begin{theorem}[Argument Principle]
Let \(\Omega\subset\C\) be a simply connected open set containing a positively oriented closed
piecewise smooth curve \(C\) and its interior. Suppose \(f\) is a meromorphic function on \(\Omega\) which has no zeros or poles on \(C\). Then
\begin{equation*}
    \frac{1}{2\pi i} \int_C \frac{f^{\prime}(z)}{f(z)}\, \dd z = N - P,
\end{equation*}
where \(N\) is the number of zeros of \(f\) inside \(C\) counting multiplicities, and \(P\) is the number of poles of \(f\) inside \(C\) counting multiplicities.
\end{theorem}
Note that for all $z_0\in\Omega$, $f(z)=(z-z_0)^m g(z)$ with $m\in \Z$ and $g(z)$ a holomorphic function which is nonzero in a neighborhood of $z_0$. The key idea is to apply the residue theorem to $\frac{f^{\prime}(z)}{f(z)}$.

\subsection{Some Important Theorems about Holomorphic Functions}
Now, we can give some important theorems about holomorphic functions.
\begin{theorem}[Rouché's Theorem]
    Let \(\Omega\subset\C\) be a connected open set containing a circle $C$ and its interior, and let \(f,g\colon\Omega\to\C\) be holomorphic functions. If \(|f(z)| > |g(z)|\) for all \(z\in C\), then \(f\) and \(f+g\) have the same number of zeros (counting multiplicities) inside \(C\).
\end{theorem}
We can prove Rouché's Theorem by applying the argument principle to $f_t(z)=f(z)+tg(z)$. We can conclude that
\[n_t\coloneq\frac{1}{2\pi i}\int_C\frac{f_t^{\prime}(z)}{f_t(z)}\,\dd z\]
is constant for $t\in [0,1]$.
\begin{theorem}[Open Mapping Theorem]
    If \(f\colon\Omega\to\C\) is a non-constant holomorphic function defined on a connected open set \(\Omega\subset\C\), then \(f\) is an open map.
\end{theorem}
The open mapping theorem can be proved using Rouché's theorem. Fix \(z_0\in\Omega\). Choose a small circle \(C\) centered at \(z_0\) and whose interior is contained in \(\Omega\) such that \(f(z)\neq f(z_0)\) for all \(z\in C\) (possible since zeros are isolated). Since \(C\) is compact, \(|f(z)-f(z_0)|\) has a positive lower bound on \(C\). Choose \(\varepsilon>0\) smaller than this bound. Then for every \(w\) with \(|w-f(z_0)|<\varepsilon\), on \(C\) we have
    \[
    |f(z_0)-w| < |f(z)-f(z_0)|.
    \]
    Write \(f(z)-w = g(z)+h(z)\), where \(g(z)=f(z)-f(z_0)\) and \(h(z)=f(z_0)-w\). 
    On that circle the inequality \(|h|<|g|\) holds, so Rouch\'e's theorem ensures that \(f-w\) has a zero inside the circle. 
    Hence \(w\in f(\Omega)\), and thus \(D(f(z_0),\varepsilon)\subset f(\Omega)\), proving that \(f(\Omega)\) is open. Applying the same argument to any open subset \(U\subset\Omega\) (\(f\) is non-constant on $U$) shows that \(f(U)\) is open; therefore \(f\) is an open map.
\begin{theorem}[Maximum Modulus Principle]
Let \(\Omega\subset\C\) be a connected open set, and let \(f\colon\Omega\to\C\) be a holomorphic function. If there exists a point \(z_0\in\Omega\) such that \(|f(z_0)| \geq |f(z)|\) for all \(z\in\Omega\), then \(f\) is constant.
\end{theorem}
By the open mapping theorem, the proof is trivial.
\begin{corollary}\label{cor:maximum modulus}
Let \(\Omega\subset\C\) be a bounded connected open set, and let \(f\) be holomorphic on $\Omega$ and continuous on $\overline{\Omega}$. Then $\sup_{z\in\overline{\Omega}} |f(z)| = \sup_{z\in\partial \Omega} |f(z)|$.
\end{corollary}
\begin{remark}
    The boundedness assumption in Corollary \ref{cor:maximum modulus} is necessary. We give a counterexample. Take $\Omega=\{z\in\C\mid \re(z)>0\}$ and $f(z)=e^z$. Then $\sup_{z\in\overline{\Omega}} |f(z)|=+\infty$, but $\sup_{z\in\partial \Omega} |f(z)| = 1$.
\end{remark}


Now we give an important lemma about holomorphic functions on the unit disk.
\begin{lemma}[Schwarz Lemma]
Let \(f\colon\mathbb{D}\to\mathbb{D}\) be a holomorphic function with \(f(0)=0\), where \(\mathbb{D}=\{z\in\C\colon|z|<1\}\) is the unit disk. Then \(|f(z)|\leq |z|\) for all \(z\in\mathbb{D}\), and \(|f^{\prime}(0)|\leq 1\). Moreover, if there exists a point \(z_0\in\mathbb{D}\setminus\{0\}\) such that \(|f(z_0)|=|z_0|\) or if \(|f^{\prime}(0)|=1\), then \(f(z)=e^{i\theta} z\) for some \(\theta\).
\end{lemma}
\begin{remark}
    By Schwarz lemma, we can compute the automorphism group of the unit disk, which is 
    \[\left\{e^{i\theta}\frac{a-z}{1-\bar{a}z}\colon a\in\mathbb{D}, \theta\in\R\right\}.\]
    Suppose $\psi_a(z) = \frac{a-z}{1-\bar{a}z}$, we have $\psi_a\circ \psi_a=id$.
\end{remark}
The key idea of the proof is to consider the function $g(z)=f(z)/{z}$ and then apply the maximum modulus principle.
\begin{theorem}[Schwarz-Pick Theorem]
Let \(f\colon\mathbb{D}\to\mathbb{D}\) be a holomorphic function. Then for all \(z_1,z_2\in\mathbb{D}\),
\begin{equation*}
    \left|\frac{f(z_1)-f(z_2)}{1-\overline{f(z_2)}f(z_1)}\right| \leq \left|\frac{z_1-z_2}{1-\overline{z_2}z_1}\right|,
\end{equation*}
and for all \(z\in\mathbb{D}\),
\begin{equation*}
    \frac{|f^{\prime}(z)|}{1-|f(z)|^2} \leq \frac{1}{1-|z|^2}.
\end{equation*}
\end{theorem}
The key idea of the proof is to consider the function $\psi_{f(z_2)}\circ f\circ \psi_{z_2}$ and then apply Schwarz lemma.

Finally, we give the existence of holomorphic primitives and holomorphic logarithms. In Section \ref{integration along curves}, we have proved that holomorphic functions on a disk have holomorphic primitives. Now we consider the general case.
\begin{theorem}[Existence of Holomorphic Primitive]\label{th:primitive}
Let \(\Omega\subset\C\) be a simply connected open set. Then all holomorphic functions \(f\colon\Omega\to\C\) have holomorphic primitives on \(\Omega\).
\end{theorem}
Fixing $z_0 \in \Omega$, we can define
\begin{equation*}
    F(z) = \int_{\gamma} f(\omega)\, \dd \omega,
\end{equation*}
where \(\gamma\) is a piecewise smooth curve in \(\Omega\) from \(z_0\) to \(z\). By the simple connectedness of \(\Omega\) and Proposition \ref{homotopic prop}, \(F(z)\) is well-defined. Then we can easily verify that \(F\) is a holomorphic primitive of \(f\).
\begin{remark}
    Simple connectedness of \(\Omega\) implies connectedness of \(\Omega\). For open subsets of $\C$, connectedness is equivalent to path connectedness. So we can define the holomorphic primitive by the integral along a path from a fixed point to \(z\).
\end{remark}
\begin{corollary}
Let \(\Omega\subset\C\) be a simply connected open set containing a closed piecewise smooth curve $\gamma$ and its interior. Then for all holomorphic functions \(f\colon\Omega\to\C\), we have
\begin{equation*}
    \int_{\gamma} f(z)\, \dd z = 0.
\end{equation*}
\end{corollary}

\begin{theorem}[Existence of Holomorphic Logarithm]\label{th:primitive log}
Let \(\Omega\subset\C\) be a simply connected open set. Then for all \(f\in\hol(\Omega)\) which have no zeros, there exists a holomorphic logarithm \(g\coloneq \log(f(z))\) such that \(g^{\prime}(z)=f'(z)/{f(z)}\) for all \(z\in\Omega\).
\end{theorem}
It is trivial by Theorem \ref{th:primitive} since $f'(z)/{f(z)}$ is holomorphic on $\Omega$.
\begin{remark}
    We can also write $f(z) = e^{g(z)}$.
\end{remark}
\begin{example}
    \begin{enumerate}
        \item  The simple connectedness assumption is necessary. We give a counterexample. Consider the function \(f(z)=z\) defined on \(\Omega=\C\setminus\{0\}\). If there exists a holomorphic logarithm \(g\colon \Omega\to\C\) such that \(g^{\prime}(z)=1/z\) for all \(z\in\Omega\), then by the residue theorem, we have
    \begin{equation*}
        \int_{C(0,1)} \frac{1}{z}\, \dd z = 2\pi i \neq 0,
    \end{equation*}
    which contradicts the existence of holomorphic primitive of \(1/z\) on \(\Omega\).
        \item In the slit plane $\C\setminus(-\infty,0]$, we have the principal branch of logarithm. 
        \begin{equation*}
            \log z = \log|z| + i \arg z,
        \end{equation*}
        where \(\arg z\in(-\pi,\pi)\).
        \item For all $\alpha\in \C$, we can define the principal branch of power function in the slit plane $\C\setminus(-\infty,0]$ by
        \begin{equation*}
            z^{\alpha} = e^{\alpha \log z}.
        \end{equation*}
    \end{enumerate}
\end{example}
\section{Conformal Maps}
Question: Given two open sets \(\Omega_1, \Omega_2\subset\C\), does there exist a biholomorphism \(f\colon\Omega_1\to\Omega_2\)? If such a map exists, can we find it?
\subsection{Examples}
Here are some examples of biholomorphic maps.
\begin{example}
    \begin{enumerate}
        \item The map \(f(z)=az+b\) with \(a\neq 0\) is a biholomorphism from \(\C\) to \(\C\).
        \item The map \(f(z)=e^{z}\) is a biholomorphism from the strip \(S=\{z\in\C\colon -\pi<\im(z)<\pi\}\) to \(\C\setminus(-\infty,0]\).
        \item The map \(f(z)=-\frac{1}{2}\left(z+1/z\right)\) is a biholomorphism from the half-disk \(\{z=x+iy\in\C\colon y>0, |z|<1\}\) to the upper half-plane $\{z\in\C\colon \im(z)>0\}$.
        \item The map \(e^{i\theta}\psi_a(z) = e^{i\theta}\frac{a-z}{1-\bar{a}z}\) with \(a\in\mathbb{D}\) and \(\theta \in \R\) is a biholomorphism from the unit disk \(\mathbb{D}\) to \(\mathbb{D}\).
        \item The map \(f(z)=\frac{z-z_0}{z-z_1}\) with \(z_0,z_1\in\C, z_0\neq z_1\) is a biholomorphism from \(\C\setminus\{z_1\}\) to \(\C\setminus\{1\}\).
        \item The map \(f(z)=\frac{az+b}{cz+d}\) with \(ad-bc\neq 0\) is a biholomorphism from \(\C_{\infty}\) to \(\C_{\infty}\), where \(\C_{\infty}=\C\sqcup\{\infty\}\) is the Riemann sphere.
    \end{enumerate}
\end{example}
Actually, $\auth(\C)=\{az+b \mid a\neq 0\}$, $\auth(\mathbb{D})=\{e^{i\theta}\psi_a \mid a\in\mathbb{D}, \theta \in \R\}$, $\auth(\C_{\infty})=\left\{\frac{az+b}{cz+d} \mid ad-bc\neq 0\right\}$. 

Consider the upper half-plane \(\mathbb{H}=\{z\in\C\colon \im(z)>0\}\). The map \(F(z)=\frac{z-i}{z+i}\) is a biholomorphism from \(\mathbb{H}\) to \(\mathbb{D}\) and the map \(G(z)=i\frac{1+z}{1-z}=F^{-1}(z)\) is a biholomorphism from \(\mathbb{D}\) to \(\mathbb{H}\). Hence \(\mathbb{H}\) and \(\mathbb{D}\) are biholomorphically equivalent.
In particular, we can conclude that \(\auth(\mathbb{H})=\{G\circ f \circ F \mid f\in \auth(\mathbb{D})\}\). And we have the following theorem.
\begin{theorem}
    Every automorphism of $\mathbb{H}$ is a fractional linear transformation taking the form $f_M$, where
    \begin{equation*}
        f_M(z) = \frac{az+b}{cz+d}, \quad M=\begin{pmatrix}
            a & b\\
            c & d
        \end{pmatrix} \in \mathrm{SL}_2(\R).
    \end{equation*}
    Moreover, $\auth(\mathbb{H})\simeq \mathrm{PSL}_2(\R)$.
\end{theorem}

\begin{proposition}\label{inj hol}
    If $f\colon \Omega_1\to \Omega_2$ is holomorphic and injective, then $f^{\prime}(z)\neq 0$ for all $z\in \Omega_1$. In particular, the inverse of $f$ defined on its range is holomorphic, and thus the inverse of a conformal map (holomorphic bijection) is also holomorphic.
\end{proposition}
The key idea of the proof is to consider the Taylor series of $f$ at a point $z_0\in \Omega_1$. If $f^{\prime}(z_0)=0$, then the Taylor series is
\begin{equation*}
    f(z) = f(z_0) + a_n (z-z_0)^n + a_{n+1} (z-z_0)^{n+1} + \cdots,
\end{equation*}
with $n\geq 2$ and $a_n\neq 0$. Then we can easily find two different points $z_1,z_2$ in a small neighborhood of $z_0$ by Rouché's theorem such that $f(z_1)=f(z_2)$, which contradicts the injectivity of $f$. 



\subsection{Riemann Mapping Theorem}
This is one of the most important theorems in complex analysis.
\begin{theorem}[Riemann Mapping Theorem]
Let \(\Omega\subset\C\) be a simply connected open set which is not equal to \(\C\). Then for all \(z_0\in\Omega\), there exists a unique conformal map \(f\colon\Omega\to\mathbb{D}\) such that \(f(z_0)=0\) and \(f^{\prime}(z_0)>0\).
\end{theorem}
Before proving the Riemann mapping theorem, we need some preparations.
\begin{theorem}[Hurwitz's Theorem]
    Suppose $(f_n)$ is a sequence of holomorphic functions defined on a connected open set $\Omega\subseteq \C$. And $(f_n)$ converges locally uniformly to a holomorphic function $f\colon \Omega\to\C$. Then
    \begin{enumerate}
        \item If for all $n\in \N$, $f_n$ is identically nonzero, then either $f$ is identically zero or $f$ is identically nonzero.
        \item If for all $n\in \N$, $f_n$ is injective, then either $f$ is constant or $f$ is injective.
    \end{enumerate} 
\end{theorem}
For the first part, the key idea of the proof is to prove that if $f$ is not identically zero, then $f_n^{\prime}/f_n$ converges uniformly to $f^{\prime}/f$ on a small circle where $f(z)$ does not vanish. Then we can apply the argument principle to conclude. 
The second part can be proved similarly by considering $g_n(z)=f_n(z)-f_n(w)$ for $w\in \Omega$.
\begin{theorem}[Montel's Theorem]
    Let \(\Omega\subset\C\) be an open set, and let \(\mathcal{F}\) be a family of holomorphic functions \(f\colon\Omega\to\C\). Then TFAE: 
    \begin{itemize}
        \item \(\mathcal{F}\) is normal, i.e., every sequence in \(\mathcal{F}\) has a subsequence which converges uniformly on every compact subset of \(\Omega\) to a holomorphic function (this limit function need not belong to $\mathcal{F}$).
        \item \(\mathcal{F}\) is uniformly bounded on every compact subset of \(\Omega\), i.e., for every compact subset \(K\subset\Omega\), there exists a constant \(M_K>0\) such that \(|f(z)|\leq M_K\) for all \(z\in K\) and all \(f\in \mathcal{F}\).
        \item \(\mathcal{F}\) is locally uniformly bounded, i.e., for every point \(z_0\in\Omega\), there exists a neighborhood \(U\subset\Omega\) of \(z_0\) and a constant \(M_{U}>0\) such that \(|f(z)|\leq M_{U}\) for all \(z\in U\) and all \(f\in \mathcal{F}\).
    \end{itemize}
\end{theorem}
\begin{proof}
    The equivalence between the second and third statements is trivial. The implication from the first to the second statement is also trivial by the definition of normal family. Now we prove the implication from the third to the first statement. 

    Suppose $(f_n)$ is a sequence in $\mathcal{F}$. We need to show that it has a subsequence converging uniformly on compact subsets of $\Omega$ to a holomorphic function.
    \begin{itemize}
    \item\textbf{Step 1: Local uniform boundedness implies equicontinuity on compact sets.}
    
    Let $K \subset \Omega$ be compact. Since $\mathcal{F}$ is locally uniformly bounded, for each $z \in K$, there exists $r_z > 0$ and $M_z > 0$ such that $D(z, 2r_z) \subset \Omega$ and $|f(w)| \leq M_z$ for all $w \in D(z, 2r_z)$ and all $f \in \mathcal{F}$.
    The open discs $\{D(z, r_z/2) \colon z \in K\}$ cover $K$, so by compactness, there exists a finite subcover $\{D(z_i, r_{z_i}/2)\}_{i=1}^N$. Let $M = \max\left\{M_{z_1}, \dots, M_{z_N}\right\}$ and $r = \min\left\{r_{z_1}, \dots, r_{z_N}\right\}$.
    
    Now take any $z_1, z_2 \in K$ with $|z_1 - z_2| < r/2$. Then $z_1 \in D(z_i, r_{z_i}/2)$ for some $i$, and $z_2 \in D(z_i, r_{z_i})$ since $|z_2 - z_i| \leq |z_2 - z_1| + |z_1 - z_i| < r/2 + r_{z_i}/2 \leq r_{z_i}$.
    By Cauchy's integral formula for derivatives, for any $f \in \mathcal{F}$ and $z \in D(z_i, r_{z_i})$, we have:
    \[
    f'(z) = \frac{1}{2\pi i} \int_{C\left(z_i, 2r_{z_i}\right)} \frac{f(\zeta)}{(\zeta - z)^2}\,\dd\zeta.
    \]
    On the circle $|\zeta - z_i| = 2r_{z_i}$, we have $|\zeta - z| \geq |\zeta - z_i| - |z_i - z| > 2r_{z_i} - r_{z_i} = r_{z_i}$, so
    \[
    |f'(z)| \leq \frac{2M}{r_{z_i}} \leq \frac{2M}{r}.
    \]
    Thus, $\mathcal{F}$ has uniformly bounded derivatives on $K$. Since $[z_1, z_2] \subset D(z_i, r_{z_i})$, we obtain
    \[
    |f(z_1) - f(z_2)| \leq \sup_{\xi \in [z_1,z_2]} |f'(\xi)| \cdot |z_1 - z_2| \leq \frac{2M}{r} |z_1 - z_2|.
    \]
    This shows $\mathcal{F}$ is equicontinuous on $K$.

    \item\textbf{Step 2: Application of Arzelà-Ascoli theorem.}
    
    Let $\{K_j\}_{j=1}^\infty$ be an exhaustion of $\Omega$ by compact sets (i.e., $K_j \subset \text{int}(K_{j+1})$ and $\bigcup_{j=1}^\infty K_j = \Omega$). The existence of such an exhaustion is trivial. 
    We consider the set $K_j=\{z\in \Omega\mid d(z,\partial \Omega)\geq 1/j,\, |z|\leq j\}$ (with the convention \(d(z,\varnothing)=+\infty\)).
    
    On $K_1$, the family $\mathcal{F}$ is uniformly bounded and equicontinuous, so by Arzelà-Ascoli, $(f_n)$ has a subsequence $(f_{n,1})$ converging uniformly on $K_1$.
    On $K_2$, the sequence $(f_{n,1})$ has a subsequence $(f_{n,2})$ converging uniformly on $K_2$.
    Continue this process to get subsequences $(f_{n,j})$ converging uniformly on $K_j$.
    Now consider the diagonal sequence $g_n = f_{n,n}$. This sequence converges uniformly on every $K_j$, hence on every compact subset of $\Omega$.

    \item\textbf{Step 3: The limit is holomorphic.}
    
    Let $g$ be the limit function. Since each $g_n$ is holomorphic and the convergence is uniform on compact sets, by Weierstrass convergence theorem, $g$ is holomorphic on $\Omega$. This completes the proof that $\mathcal{F}$ is normal.
    \end{itemize}
\end{proof}
\begin{proof}[Proof of Riemann Mapping Theorem]
    Once we have established the technical results above, the rest of the proof is very elegant. It consists of four steps, which we isolate.
    \begin{itemize}
        \item Step 1: Uniqueness. Suppose \(f_1,f_2\colon\Omega\to\mathbb{D}\) are two conformal maps such that \(f_1(z_0)=f_2(z_0)=0\) and \(f_1^{\prime}(z_0), f_2^{\prime}(z_0)>0\). Consider the map
        \begin{equation*}
            g(z) = f_2 \circ f_1^{-1}(z).
        \end{equation*}
        Then \(g\colon\mathbb{D}\to\mathbb{D}\) is a conformal map with \(g(0)=0\). By the Schwarz lemma, we have $g(z)=e^{i\theta}z$. Since \(g^{\prime}(0) = f_2^{\prime}(z_0)/{f_1^{\prime}(z_0)}>0\), we have \(\theta=0\). Hence \(g(z)=z\) and \(f_1=f_2\).
        \item Step 2: We claim that $\Omega$ is  conformally equivalent to an open subset of the unit disk that contains the origin. Choose a point \(a\in\C\setminus\Omega\) (possible since \(\Omega\neq \C\)). Since $z-a\neq 0$ on $\Omega$,
        by Theorem \ref{th:primitive log}, there exists a holomorphic logarithm \(h\colon\Omega\to\C\) such that \(e^{h(z)}=z-a\) for all \(z\in\Omega\), which proves in particular that $h$ is injective. 
        Hence, we have 
        \begin{equation*}
            h(z)\neq h(z_0)+2\pi i,\,\forall z\in\Omega.
        \end{equation*}
        Since \(h(\Omega)\) is open (by the open mapping theorem), there exists \(r>0\) such that \(D(h(z_0),r)\subset h(\Omega)\). Because \(e^{h(z)}=z-a\) is injective, the shifted disk \(D(h(z_0)+2\pi i,r)\) is disjoint from \(h(\Omega)\). Indeed, if \(w\in D(h(z_0),r)\) and \(w+2\pi i=h(z_1)\) for some \(z_1\in\Omega\), then \(e^{h(z_1)}=e^{w}=e^{h(z_w)}\) for the unique \(z_w\in\Omega\) with \(h(z_w)=w\), 
        contradicting the injectivity of \(h\). Thus this shifted disk contains no points of the image. 
        Now we can define a holomorphic function $H$ on $\Omega$ by
        \begin{equation*}
            H(z) = \frac{1}{h(z)-(h(z_0)+2\pi i)}.
        \end{equation*} 
        Since $h$ is injective, by Proposition \ref{inj hol}, $H\colon \Omega \to H(\Omega)$ is a conformal map. Moreover, $H(\Omega)$ is bounded by $1/r$. We may therefore translate and rescale the function $H$ in order to obtain a conformal map from $\Omega$ to an open subset of $\mathbb{D}$ that contains the origin.
        \item Step 3: By the second step, we may assume that \(\Omega\)  is an open subset of \(\mathbb{D}\) with \(0\in\Omega\). Consider the family
        \begin{equation*}
            \mathcal{F} = \{f\colon\Omega\to\mathbb{D}\mid f\text{ is holomorphic, injective}, f(0)=0\}.
        \end{equation*}
        $\mathcal{F}$ is non-empty since the identity map is in $\mathcal{F}$. Moreover, $\mathcal{F}$ is uniformly bounded by definition. 
        
        \textbf{Goal}: find \(f\in\mathcal{F}\) such that \(|f^{\prime}(0)|\) is maximal.

        By Cauchy's estimates, $|f^{\prime}(0)|$ is bounded. Suppose $s=\sup_{f\in \mathcal{F}}|f^{\prime}(0)|$. Then we can choose a sequence $(f_n)$ in $\mathcal{F}$ such that $\lim_{n\to \infty} |f_n^{\prime}(0)|=s$. 
        By Montel's theorem, there exists a subsequence $(f_{n_k})$ which converges uniformly on every compact subset of $\Omega$ to a holomorphic function $f\colon \Omega \to \overline{\mathbb{D}}$. 
        Since $s\geq 1$ (because the identity map is in $\mathcal{F}$), $f$ is non-constant. Since \(f\) is non-constant and \(|f|\le 1\) on \(\Omega\), the maximum modulus principle implies that \(|f(z)|<1\) for all \(z\in\Omega\). Hence \(f\) indeed maps into \(\mathbb D\). Moreover, by Hurwitz's theorem, $f$ is injective.
        Since we clearly have $f(0)=0$, we obtain $f\in \mathcal{F}$ with $|f^{\prime}(0)|=s$.
        \item Step 4: We claim that $f$ is a conformal map from $\Omega$ to $\mathbb{D}$. Since $f$ is already injective, it suffices to prove \(f(\Omega)=\mathbb{D}\). 
        If this were not true, we could construct a function in $\mathcal{F}$ with derivative at $0$ greater than $s$. Suppose there exists a point \(\alpha\in\mathbb{D}\setminus f(\Omega)\). Consider the map 
        \[\psi_\alpha=\frac{\alpha-z}{1-\overline{\alpha}z},\]
        which is an automorphism of \(\mathbb{D}\) that interchanges \(\alpha\) and \(0\). Since $\Omega$ is simply connected, so is $U=(\psi_\alpha\circ f)(\Omega)$, and moreover, $0\notin U$. 
        It is therefore possible to define a square root function on $U$ by
        \begin{equation*}
            u(w) = e^{\frac{1}{2}\log(w)},\, w\in U.
        \end{equation*}
         Next, consider the function
         \begin{equation*}
            F=\psi_{u(\alpha)}\circ u \circ \psi_\alpha \circ f.
         \end{equation*}
         We claim that $F\in\mathcal{F}$. Indeed, $F$ is holomorphic, maps $0$ to $0$ and maps into the unit disk (since each factor does). Moreover, $F$ is injective. 

         Denote by $\varphi$ the square function $\varphi(z)=z^2$. Then we have
         \begin{equation*}
            f=\psi_\alpha^{-1}\circ \varphi \circ \psi_{u(\alpha)}^{-1} \circ F\eqcolon \Phi\circ F.
         \end{equation*}
         \(\Phi\) maps \(\mathbb{D}\) into \(\mathbb{D}\) with \(\Phi(0) = 0\), but it is not injective because \(\varphi\) is not. By the last part of the Schwarz lemma, we conclude that \(|\Phi'(0)| < 1\). The proof is complete once we observe that
         \[f'(0) = \Phi'(0) F'(0),\]
         and thus \[|f'(0)| < |F'(0)|,\] contradicting the maximality of \(|f'(0)|\) in \(\mathcal{F}\). Finally, we multiply \(f\) by a complex number of absolute value 1 so that \(f'(0) > 0\), which ends the proof.
    \end{itemize}
\end{proof}
This proof is adapted from \textit{Complex Analysis} by Elias M. Stein and Rami Shakarchi.


\section{Some Interesting Problems}
\begin{problem}
    Although $\Omega=\C\setminus\{0\}$ is not simply connected, prove that $e^{\frac{1}{z}}-\frac{1}{z}$ has a holomorphic primitive on $\Omega$.
\end{problem}
\begin{solution}
    Consider the Laurent series of $e^{\frac{1}{z}}-\frac{1}{z}$ at $0$, we have
    \begin{equation*}
        e^{\frac{1}{z}}-\frac{1}{z} = \sum_{n=0}^{+\infty} \frac{1}{n!} \cdot \frac{1}{z^n} - \frac{1}{z} = 1 + \frac{1}{2!}\cdot \frac{1}{z^2} + \frac{1}{3!}\cdot \frac{1}{z^3} + \cdots.
    \end{equation*}
    So we can see that $e^{\frac{1}{z}}-\frac{1}{z}$ has a primitive
    \begin{equation*}
        F(z) = z - \sum_{n=2}^{+\infty} \frac{1}{n!} \cdot \frac{1}{(n-1) z^{n-1}}.
    \end{equation*}
    Since the series is convergent absolutely on every compact subset of $\Omega$, by Weierstrass convergence theorem, $F$ is holomorphic on $\Omega$.
\end{solution}
\begin{problem}
    Suppose $\Omega=\C\setminus[0,1]$. Prove that there exist two distinct well-defined holomorphic functions on $\Omega$ such that their squares are $z\mapsto z(z-1)$.
\end{problem}
\begin{solution} We can prove it in 3 steps.
    \begin{itemize}
        \item Step 1: If we can find two distinct holomorphic functions $h_1$ and $h_2$ satisfying $h_1^2(z)=h_2^2(z)=z(z-1)$, then $h_1=-h_2$.
        \item Step 2: On $\C\setminus (-\infty,1]$, we can define 
        \begin{equation*}
            h_1(z) = \exp\left({\frac{1}{2}\log z + \frac{1}{2}\log (z-1)}\right).
        \end{equation*}
        \item Step 3: We can extend $h_1$ to $\Omega$ by defining $h_1(x)=-\sqrt{x(x-1)}$ for $x\in \R_{<0}$. Then it suffices to verify that the extension is holomorphic at every point in $(-\infty,0)$.
    \end{itemize}
\end{solution}
\begin{problem}
    For which $\alpha\in\C$, does $f_\alpha\colon \C\setminus (-\infty,0]\to \C$ defined by $f_\alpha(z)=z^\alpha$ have an analytic continuation to $\C^{*}$ or $\C$?
\end{problem}
\begin{solution}
    If $f_\alpha$ has an analytic continuation on $\C^{*}$, we have 
    \begin{equation*}
        \lim_{\theta\rightarrow\pi}f_\alpha(e^{i\theta})= \lim_{\theta\rightarrow \pi}f_\alpha(e^{-i\theta}),
    \end{equation*}
    which gives us $e^{2\pi i\alpha}=1$. Hence $\alpha\in\Z$. For all $\alpha\in \Z$, $f_\alpha$ can be extended to $\C^{*}$ by defining $f_\alpha(x)=x^{\alpha}$ for all $x\in (-\infty,0]$.

    If $f_\alpha$ has an analytic continuation on $\C$, we obtain $\alpha\in\N$ since $0$ is a pole if $\alpha<0$. For all $\alpha\in \Z_{>0}$, we can extend $f_\alpha$ to $\C$ by defining $f_\alpha(0)=0$. And for $\alpha=0$, we can extend $f_0$ to $\C$ by defining $f_0(0)=1$.
\end{solution}
\begin{problem}
    Suppose $X=\{z\in \C\mid \re(z)=0,|\im(z)|\geq 1\}$ and $\Omega=\C\setminus X$.
    \begin{enumerate}
        \item[(a)] Prove that there exists a unique holomorphic function $f\colon \Omega\to \C$ such that $f^{\prime}(z)=\frac{1}{1+z^2}$ and $f(0)=0$.
        \item[(b)] Does $z\mapsto \frac{1}{1+z^2}$ have a holomorphic primitive on $\C\setminus \{\pm i\}$?
        \item[(c)] Does $z\mapsto \frac{1}{1+z^2}$ have a holomorphic primitive on $\C\setminus [-i,i]$?
    \end{enumerate} 
\end{problem}
\begin{solution}
    \begin{enumerate}
        \item [(a)] Since $\Omega$ is simply connected, by Theorem \ref{th:primitive}, there exists a holomorphic primitive $f$ of $z\mapsto \frac{1}{1+z^2}$ on $\Omega$. And the uniqueness is trivial.
        \item [(b)] No. If there exists a holomorphic primitive $F$ of $z\mapsto \frac{1}{1+z^2}$ on $\C\setminus \{\pm i\}$, by the residue theorem, we have
        \begin{equation*}
            \int_{C(i,1)} \frac{1}{1+z^2}\, \dd z = \pi \neq 0,
        \end{equation*}
        which contradicts the existence of holomorphic primitive.
        \item [(c)] Yes. First, we can define a primitive $F$ of $z\mapsto \frac{1}{1+z^2}$ on $\C\setminus (-i\infty,i]$ by
        \begin{equation*}
            F(z) = \int_{[2i,z]} \frac{1}{1+\omega^2}\, \dd \omega,
        \end{equation*}
        Then we can extend $F$ to $\C\setminus [-i,i]$ by defining 
        \begin{equation*}
            F(x) = \int_{\gamma_x} \frac{1}{1+\omega^2}\, \dd \omega,\, x\in (-i\infty,-i),
        \end{equation*}
        where $\gamma_x$ is a piecewise smooth path in $\C\setminus [-i,i]$ from $2i$ to $x$ which is contained in either the right half-plane or the left half-plane. 
        Since
        \[\res_{z=i} \frac{1}{1+z^2} + \res_{z=-i} \frac{1}{1+z^2} = 0,\]
        by the residue theorem, for all simple closed piecewise smooth curve $\gamma$ in $\C\setminus [-i,i]$ around the segment $[-i,i]$, the integral along $\gamma$ is $0$.
        Hence, the extension is well-defined. Finally, we only need to verify that the extension is holomorphic at every point in $(-i\infty,-i)$. It is trivial.
    \end{enumerate}
\end{solution}


\begin{problem}
    Prove that for all holomorphic function $f\colon \C\to \C$, there is a holomorphic function $g\colon \C\to \C$ such that $f(z)=g(z+1)-g(z)$.
\end{problem}
\begin{solution}
    Let  
\[
f(z)=\sum_{k=0}^{\infty} a_k z^k
\]
be the Taylor expansion of the entire function \(f\). For each \(k\ge 0\), let \(r_k\) be the unique odd multiple of \(\pi\) in the interval \((k,k+2\pi)\), and define
\[
F_k(z)=\frac{k!}{2\pi i}\oint_{|\omega|=r_k}\frac{e^{z\omega}}{\omega^{k+1}(e^\omega-1)}\,\dd\omega .
\]

First note that the contour avoids the zeros of \(e^\omega-1\), since those zeros have modulus \(2\pi n\), while \(r_k\) is an odd multiple of \(\pi\). Hence each \(F_k\) is well-defined and, by standard differentiation under the integral sign, is entire.

Now compute:
\[
F_k(z+1)-F_k(z)
= \frac{k!}{2\pi i}\oint_{|\omega|=r_k}
\frac{e^{z\omega}(e^\omega-1)}{\omega^{k+1}(e^\omega-1)}\,\dd\omega
= \frac{k!}{2\pi i}\oint_{|\omega|=r_k}\frac{e^{z\omega}}{\omega^{k+1}}\,\dd\omega .
\]
By Cauchy's integral formula,
\[
\oint_{|\omega|=r_k}\frac{e^{z\omega}}{\omega^{k+1}}\,\dd\omega
= \frac{2\pi i}{k!} z^k,
\]
so
\[
F_k(z+1)-F_k(z)=z^k.
\]

We need an estimate for \(F_k\). Write \(\omega=x+iy\) with \(x^2+y^2=r_k^2\). Since \(r_k\) is an odd multiple of \(\pi\), we claim that on this circle
\[
|e^\omega-1|\ge \frac12.
\]
Indeed, suppose \(|e^\omega-1|<1/2\). Then
\[
e^x=|e^\omega|\ge 1-|e^\omega-1|>\frac12,\qquad
e^x\le 1+|e^\omega-1|<\frac32,
\]
so \(|x|<1\). Also,
\[
|e^\omega-1|^2=(e^x\cos y-1)^2+(e^x\sin y)^2<\frac14,
\]
which implies \(e^x\cos y>1/2\), hence \(\cos y>0\). Thus \(y\) lies in one of the intervals \((2\pi n-\frac{\pi}{2},\,2\pi n+\frac{\pi}{2})\). If \(n=0\), then \(|y|<\pi/2\), contradicting \(|y|=\sqrt{r_k^2-x^2}\ge \sqrt{\pi^2-1}>\pi/2\). If \(|n|\ge 1\), then \(|y|\ge 2\pi|n|-\pi/2\geq 3\pi/2\). By symmetry, we may assume $y>0$. A short check shows that this forces
\[
x^2=r_k^2-y^2=(r_k+y)(r_k-y)\geq 2y\cdot \frac{\pi}{2}\ge \frac{3\pi^2}{2},
\]
contradicting \(|x|<1\). Hence \(|e^\omega-1|\ge 1/2\) on the contour.

Moreover, we have \(r_k>k\), so \(r_k^k\ge k^k\ge k!\). Therefore,
\[
|F_k(z)|
\le \frac{k!}{2\pi}\cdot 2\pi r_k\cdot
\frac{e^{|z|r_k}}{r_k^{k+1}\cdot (1/2)}
= \frac{2k!\,e^{|z|r_k}}{r_k^k}
\le 2e^{|z|r_k}
\le 2e^{(k+2\pi)|z|}.
\]

Now define
\[
g(z)=\sum_{k=0}^{+\infty} a_k F_k(z).
\]
For any compact set \(|z|\le R\), we have 
\[
|a_k F_k(z)|\le 2e^{2\pi R}\, |a_k| e^{kR}.
\]
Since \(f\) is entire, \(\limsup_{k\to\infty}|a_k|^{1/k}=0\). Thus \(\sum |a_k| e^{kR}\) converges, and so the series defining \(g\) converges uniformly on compact subsets of \(\mathbb C\). Hence \(g\) is entire.

Finally, using the fact that \(F_k(z+1)-F_k(z)=z^k\) and uniform convergence,
\[
g(z+1)-g(z)=\sum_{k=0}^{\infty} a_k\bigl(F_k(z+1)-F_k(z)\bigr)=\sum_{k=0}^{\infty} a_k z^k= f(z).
\]
Thus \(g\) is an entire function satisfying \(f(z)=g(z+1)-g(z)\), as required.
\end{solution}


\begin{problem}
    Prove that there is no entire function $f$ such that $f\circ f(z)=e^z$ for all $z\in\C$.
\end{problem}
\begin{solution}
    Assume, for contradiction, that there exists an entire function $f$ such that
\[
f(f(z))=e^z \quad \forall z\in\mathbb C.
\]

We first show that $f$ has no zeros. Suppose $f(a)=0$ for some $a\in\mathbb C$. Then $a\notin f(\mathbb C)$; indeed, if $f(z_0)=a$ for some $z_0$, then
\[
f(f(z_0))=f(a)=0,
\]
but $f(f(z_0))=e^{z_0}\neq 0$, a contradiction.
Since $\C\setminus\{0\}=f\circ f(\C)\subseteq f(\C)$, we must have $a=0$. Therefore, $f(0)=0$. However, this implies that $f(f(0))=f(0)=0$, which contradicts the fact that $f(f(0))=e^0=1$. Hence $f$ has no zeros.

Since $f$ is entire and nowhere zero, by Theorem \ref{th:primitive log}, there exists an entire function $g$ such that
\[
f(z)=e^{g(z)}.
\]
Substituting into the original equation gives
\[
e^z=f(f(z))=e^{g(f(z))}.
\]
Hence
\[
g(f(z))=z+2\pi i\cdot  k(z)
\]
for some integer-valued continuous function $k(z)$. Since $\mathbb C$ is connected, $k\colon \C\to \Z$ is constant; thus
\[
g(f(z))=z+C
\]
for some constant $C$.

Now we prove that $f$ is injective. If $f(z_1)=f(z_2)$, then applying $g$ to both sides yields
\[
z_1+C=g(f(z_1))=g(f(z_2))=z_2+C,
\]
so $z_1=z_2$. Hence $f$ is injective.

Then the composition $f\circ f$ is also injective. However, $e^z$ is not injective. This contradicts the identity $f(f(z))=e^z$. Therefore, no such entire function $f$ exists. \qedsymbol
\end{solution}

\begin{problem}
    Suppose $f$ is an entire function such that $|f(z)|=1$ for all $z$ with $|z|=1$. Prove that $f(z)=az^n$ with $|a|=1$ and $n\in \N$.
\end{problem}
\begin{solution}
    Define the function 
    \[
    h(z)=f(z)\overline{f(1/\bar{z})}, \quad z\in \mathbb C\setminus\{0\}.
    \]
    Then \(h\) is holomorphic on \(\mathbb C\setminus\{0\}\). For \(|z|=1\), we have \(1/\bar{z}=z\), and since \(|f(z)|=1\),
    \[
    h(z)=f(z)\overline{f(z)}=|f(z)|^2=1.
    \]
    By the identity theorem, \(h(z)\equiv 1\) on \(\mathbb C\setminus\{0\}\). Hence \(f(z)\neq 0\) for all \(z\neq 0\), and 
    \[
    f(z)=\overline{f(1/\bar{z})}^{-1}\eqqcolon g(z),\,\forall z\in \C.
    \]

    If \(f(0)\neq 0\), then \(g\) extends to \(\C\) by setting \(g(0)=f(0)\). Moreover, as \(z\to\infty\), we have \(1/\bar{z}\to 0\), so 
    \[
    g(z)=\overline{f(1/\bar{z})}^{-1}\to \overline{f(0)}^{-1}.
    \]
    Thus \(g\) is bounded on \(\mathbb C\), hence constant by Liouville's theorem. So \(f\equiv a\) with \(|a|=1\).

    If \(f(0)=0\), since \(f\not\equiv 0\), write \(f(z)=z^n h(z)\) with \(n\in\mathbb N\) and \(h(0)\neq 0\). For \(|z|=1\), we have \(|h(z)|=|f(z)|=1\). By the previous case applied to \(h\), \(h\) is constant. Hence \(f(z)=a z^n\) with \(|a|=1\). \qedsymbol
\end{solution}

\begin{problem}
    Suppose $(f_n)$ is a sequence of holomorphic functions defined on an open set $\Omega\subseteq \C$. If $(f_n)$ converges pointwise to a function $f\colon \Omega\to\C$ and if for every compact $K\subset \Omega$, there exists a constant $C(K)$ such that 
    \[\sup_{z\in K}|f_n(z)|\leq C(K),\,\forall n\in\N.\] 
    Prove that $f$ is holomorphic on $\Omega$ and for all $k\in \N$ and for every compact $K\subset \Omega$, there exists a constant $C(K,k)$ such that
    \[\sup_{z\in K}\left|f_n^{(k)}(z)\right|\leq C(K,k),\,\forall n\in\N.\] 
\end{problem}
\begin{solution}
    We see that $(f_n)$ is locally uniformly bounded. By Montel's theorem, there exists a subsequence $(f_{n_k})$ which converges uniformly on every compact subset of $\Omega$ to a holomorphic function $g\colon \Omega\to\C$. Since $(f_n)$ converges pointwise to $f$, we have $f=g$ on $\Omega$. Hence $f$ is holomorphic on $\Omega$. 
    
    Let $K \subset \Omega$ be compact. Since $(f_n)$ is locally uniformly bounded, for each $z \in K$, there exists $r_z > 0$ and $M_z > 0$ such that $D(z, 2r_z) \subset \Omega$ and $|f_n(w)| \leq M_z$ for all $w \in D(z, 2r_z)$ and all $n \in \mathbb{N}$.
    The open discs $\{D(z, r_z) \colon z \in K\}$ cover $K$, so by compactness, there exists a finite subcover $\{D(z_i, r_{z_i})\}_{i=1}^N$. Let $M = \max\left\{M_{z_1}, \dots, M_{z_N}\right\}$ and $r = \min\left\{r_{z_1}, \dots, r_{z_N}\right\}$.
    
    By Cauchy's integral formula for derivatives, for any $n\in \N$ and $z \in K$:
    \[
    f_n^{(k)}(z) = \frac{k!}{2\pi i} \int_{C\left(z_i, 2r_{z_i}\right)} \frac{f_n(\omega)}{(\omega - z)^{k+1}} \dd\omega.
    \]
    On the circle $|\omega - z_i| = 2r_{z_i}$, we have $|\omega - z| \geq |\omega - z_i| - |z_i - z| > 2r_{z_i} - r_{z_i} = r_{z_i}$, so
    \[
    \left|f_n^{(k)}(z)\right| \leq \frac{k! \cdot 2M}{r_{z_i}^{k}} \leq \frac{k! \cdot 2M}{r^{k}}.
    \]

    Hence we can take $C(K,k)=k! \cdot 2M/{r^{k}}$.
\end{solution}
\begin{problem}
    If $\Omega$ is a connected open bounded set in $\C$ with $C^1$ boundary. Suppose that there exist a sequence of automorphisms $f_n\colon \Omega \to \Omega$ with a point $a\in \Omega$ such that $f_n(a)\to b\in \partial \Omega$ as $n\to +\infty$. 
    Prove that $\Omega$ is conformally equivalent to the unit disc.
\end{problem}
\begin{solution}
    By Riemann mapping theorem, it suffices to prove that $\Omega$ is simply connected. We conclude it by three steps.
    \begin{itemize}
        \item Step 1: Since $\Omega$ is bounded, by applying Montel's theorem to $g_n(z)=f_n(z)-b$, there exists a subsequence $(g_{n_k})$ which converges uniformly on every compact subset of $\Omega$ to a holomorphic function $g \colon \Omega \to \C$. Clearly, we have $g(a)=0$.
        Since $g_{n_k}$ does not have zeros, by Hurwitz's theorem, $g(z)=0$ for all $z\in \Omega$. Hence $(f_{n_k})$ converge locally uniformly to the constant function $b$.
        \item Step 2: Suppose $\gamma$ is a path in $\Omega$. We claim that for all $\varepsilon>0$, there exists $N\in\N$ such that for all $k>N$, $f_{n_k}(\im(\gamma))\subset \Omega\cap D(b,\varepsilon)$. It is trivial by the local uniform convergence of $(f_{n_k})$ to $b$.
        \item Step 3: By contradiction, suppose $\Omega$ is not simply connected. Then there exists a non-contractible closed path $\gamma$ in $\Omega$. By the above step, for all $\varepsilon>0$, there exists $N\in\N$ such that for all $k>N$, $f_{n_k}(\im(\gamma))\subset \Omega\cap D(b,\varepsilon)$. 
        Since $\partial \Omega$ is $C^1$, we can choose $\varepsilon$ small enough such that $\Omega\cap D(b,\varepsilon)$ is simply connected (This is because we can choose $\varepsilon$ small enough such that $\Omega\cap D(b,\varepsilon)$ is homeomorphic to a half disc). Hence $f_{n_k}(\im(\gamma))$ is contractible in $\Omega\cap D(b,\varepsilon)$, which contradicts the fact that $f_{n_k}$ is an automorphism of $\Omega$ since $\gamma$ is non-contractible. \qedsymbol
    \end{itemize} 
\end{solution}

\begin{problem}
    Let $r_1>r_2>1$ and $A_1=\{z\in \C\mid 1<|z|<r_1\}$, $A_2=\{z\in \C\mid 1<|z|<r_2\}$. Prove that there is no biholomorphism $f\colon A_1\to A_2$.
\end{problem}
\begin{solution}
    Suppose there exists a biholomorphism $f\colon A_1\to A_2$. Consider $u(z)=\frac{\log |z|}{\log r_2}$. Define the Dirichlet energy of $u$ on $A_2$ by 
    \[E_{A_2}(u)=\int_{A_2}|\nabla u|^2 \, \dd\sigma=\int_{0}^{2\pi}\int_{1}^{r_2} \left|\nabla \left(\frac{\log(re^{i\theta})}{\log r_2}\right)\right|^2 r \,\dd r \,\dd\theta=\frac{2\pi}{\log r_2}.\]
    Since $f$ is biholomorphic, we have \(E_{A_2}(u)=E_{A_1}(u\circ f)\). Suppose $v=u\circ f$. Since $\Delta u=0$, we have
    \[\Delta v=\Delta (u\circ f)=|f^{\prime}|^2 (\Delta u)\circ f=0.\]
    Therefore, $v$ is harmonic on $A_1$. \\[2mm]
    \textbf{Claim:} $\lim_{|z|\to r_1^-}|f(z)|=r_2,\,\lim_{|z|\to 1^+}|f(z)|=1$ or $\lim_{|z|\to r_1^-}|f(z)|=1,\,\lim_{|z|\to 1^+}|f(z)|=r_2$.\\[2mm]
    \noindent\textbf{Proof of the claim.} Since $f$ is a homeomorphism, it is proper: for every compact $K\subset A_2$, $f^{-1}(K)$ is compact in $A_1$. Fix $\rho\in(1,r_2)$ and set $C_\rho=\{w\in A_2: |w|=\rho\}$. Then $f^{-1}(C_\rho)$ is compact in $A_1$, hence has positive distance from the boundary of $A_1$.

Take a sequence $\{z_n\}\subset A_1$ with $|z_n|\to 1^+$. The values $|f(z_n)|$ lie in $(1,r_2)$. If a subsequence converged to some $\rho\in(1,r_2)$, then there exists a subsequence $f(z_{n_k})$ that would converge to a point $w_0$ with $|w_0|=\rho$, and by continuity of $f^{-1}$, $z_{n_k}\to f^{-1}(w_0)\in A_1$, contradicting $|z_n|\to 1$. Thus all accumulation points of $|f(z_n)|$ are in $\{1,r_2\}$.

    We show that only one accumulation point occurs. If not, there are subsequences $z_n', z_n''$ with $|z_n'|,|z_n''|\to 1$ but $|f(z_n')|\to 1$ and $|f(z_n'')|\to r_2$. Choose $\rho\in(1,r_2)$. For large $n$, $|f(z_n')|<\rho<|f(z_n'')|$. Connect $z_n'$ and $z_n''$ by a path $\gamma_n\subset A_1$ lying arbitrarily close to the inner boundary, i.e. $\max_{z\in \gamma_n}|z|\to 1$. The image path $f(\gamma_n)$ crosses the circle $|w|=\rho$, so there exists $\xi_n\in \gamma_n$ with $|f(\xi_n)|=\rho$. Hence $\xi_n\in f^{-1}(C_\rho)$ and $|\xi_n|\to 1$ since $|\xi_n|\leq \max_{z\in \gamma_n}|z|$, contradicting the compactness of $f^{-1}(C_\rho)$. Therefore the limit exists and is either $1$ or $r_2$.

The same argument applied to sequences with $|z_n|\to r_1^-$ shows that $|f(z_n)|$ converges to a single value in $\{1,r_2\}$. Since $f$ is bijective, the previous argument applied to $f^{-1}$ shows that the two boundary components cannot both go to the same component. Hence the two possible pairings occur. The claim is proved. $\blacksquare$
    
    According to the claim, we conclude that 
    \[ \lim_{|z|\to r_1^-}v(z)=1,\,\lim_{|z|\to 1^+}v(z)=0 \text{ or } \lim_{|z|\to r_1^-}v(z)=0,\,\lim_{|z|\to 1^+}v(z)=1.\]
    Since $v$ is harmonic on $A_1$, the boundary limits determine it uniquely (by the maximum principle applied to the difference of two such functions on compact sub-annuli). Therefore
    \[v(z)=\frac{\log |z|}{\log r_1}\text{ or }v(z)=\left(1-\frac{\log |z|}{\log r_1}\right).\]
    Hence we have
    \[E_{A_1}(v)=\int_{A_1}|\nabla v|^2 \, \dd\sigma=\frac{2\pi}{\log r_1}.\]
    However, since $r_1>r_2$, we have $E_{A_1}(v)=2\pi/\log r_1<{2\pi}/{\log r_2}=E_{A_2}(u)$, which contradicts the fact that $E_{A_2}(u)=E_{A_1}(v)$. Hence there is no biholomorphism $f\colon A_1\to A_2$. \qedsymbol
\end{solution}

\end{document}